% hookes: Hooke's law, F = -kx, always restoring (Ch 8, fig:hookes)
% Left: a block on a smooth floor tied to a wall by a spring, at natural length,
% stretched (x = +0.7 cm) and compressed (x = -0.6 cm).  The dashed guide marks
% the natural-length position of the block's face; x is measured from it.  The
% spring force on the block (crimson) points back towards that guide in both cases,
% with length proportional to |x| (1.2 cm per cm of deformation).
% Right: the graph of F = -kx; the two marked points are the two deformed states.
% Every spring has the same number of turns (8), so a stretched spring has wider
% turns and a compressed one closer turns.
\begin{tikzpicture}[x=1cm,y=1cm]
  % a coil spring from (#1,#2) to (#3,#2): straight leads of 0.15 cm, then N = 8 turns
  % drawn as the side view of a helix, x = p t/2pi + q(1 - cos t), y = A sin t
  \def\spring(#1,#2,#3){%
    \pgfmathsetmacro\p{((#3)-(#1)-0.3)/8}%
    \draw[cInk,thin line] (#1,#2) -- ++(0.15,0)
      plot[domain=0:16*pi,samples=320,smooth,variable=\t]
        ({#1+0.15+\p*\t/(2*pi)+0.09*(1-cos(\t r))},{#2+0.16*sin(\t r)})
      -- (#3,#2);}
  \def\Ln{2.2}                % natural length
  \def\bw{0.65}\def\bh{0.75}  % block
  \def\ys{0.25}               % height of the spring
  \def\yf{0.58}               % height of the force arrow
  % three rows: label / deformation x / vertical position
  \foreach \lab/\x/\y in {{natural length, $F=0$}/0/0, {stretched, $x>0$}/0.7/-1.45, {compressed, $x<0$}/-0.6/-2.9}{
    \begin{scope}[yshift=\y cm]
      \pgfmathsetmacro\L{\Ln+\x}
      % wall and floor
      \path[ground,draw=none] (-0.16,0) rectangle (0,1.05);
      \draw[outline] (0,0) -- (0,1.05);
      \path[ground,draw=none] (0,-0.14) rectangle (3.9,0);
      \draw[outline] (0,0) -- (3.9,0);
      % spring and block
      \spring(0,\ys,\L)
      \draw[body] (\L,0) rectangle (\L+\bw,\bh);
      \node[anchor=west,font=\footnotesize,text=cSteel] at (4.0,0.38) {\lab};
      % force and deformation
      \ifdim\x pt>0pt
        \draw[force] (\L,\yf) -- ++(-1.2*\x,0) node[anchor=south,pos=0.3,inner sep=2pt] {$F$};
        \draw[dim] (\Ln,-0.36) -- node[below=-1pt,cSteel] {$x$} (\L,-0.36);
      \fi
      \ifdim\x pt<0pt
        \draw[force] (\L+\bw,\yf) -- ++(-1.2*\x,0) node[anchor=west,inner sep=2pt] {$F$};
        \draw[dim] (\L,-0.36) -- node[below=-1pt,cSteel] {$|x|$} (\Ln,-0.36);
      \fi
    \end{scope}
  }
  % the natural-length position, common to all three rows
  \draw[guide] (\Ln,1.35) -- (\Ln,-3.45);
  \node[anchor=south,font=\footnotesize,text=cSteel] at (\Ln,1.35) {$x=0$};
  % the graph
  \begin{scope}[xshift=9.3cm,yshift=-1.35cm]
  \begin{axis}[bookplot, anchor=center, at={(0,0)},
      width=5.0cm, height=4.6cm, scale only axis,
      axis lines=middle, xmin=-1.15, xmax=1.15, ymin=-1.15, ymax=1.15,
      xlabel={$x$}, ylabel={$F$},
      xlabel style={anchor=north east}, ylabel style={anchor=north west,xshift=1pt}]
    \addplot[cCrim,domain=-1:1,samples=2] {-x};
    \node[cCrim,anchor=south west,inner sep=1pt] at (axis cs:-0.98,1.0) {$F=-kx$};
    \node[cSteel,anchor=north east] at (axis cs:-0.12,-0.16) {slope $-k$};
    % the two deformed states
    \draw[guide] (axis cs:0.7,0) -- (axis cs:0.7,-0.7) -- (axis cs:0,-0.7);
    \draw[guide] (axis cs:-0.6,0) -- (axis cs:-0.6,0.6) -- (axis cs:0,0.6);
    \node[point mass] at (axis cs:0.7,-0.7) {};
    \node[point mass] at (axis cs:-0.6,0.6) {};
    \node[anchor=south west,font=\footnotesize,text=cSteel,align=left] at (axis cs:0.74,-0.66)
      {stretched:\\ $F<0$};
    \node[anchor=north east,font=\footnotesize,text=cSteel,align=right] at (axis cs:-0.64,0.56)
      {compressed:\\ $F>0$};
  \end{axis}
  \end{scope}
\end{tikzpicture}
